\chapter{Teoretické základy a terminologie}
\label{chap:teorie}

Tato kapitola představuje základní pojmy a teoretická východiska, o která se opírá návrh vysvětlitelného a ovladatelného doporučovacího systému. Nejprve jsou představeny tradiční i moderní přístupy k doporučování, včetně jejich matematické formalizace. Následně je rozebrán koncept řídkých autoenkodérů (SAE) a specifika domény bodů zájmu (POI). V závěru jsou definovány principy vysvětlitelnosti a interaktivního řízení preferencí (steering).

\section{Doporučovací systémy a doména POI}
\label{sec:rs_zaklady}

Doporučovací systémy (\textit{Recommender Systems}, RS) jsou softwarové aplikace určené k~personalizovanému navrhování položek či aktivit na základě preferencí, chování a~kontextu uživatele \cite{RicciRokachShapira2022}. Jejich cílem je pomoci uživateli zorientovat se v~množství dostupných možností a~zvýšit pravděpodobnost, že si vybere relevantní položku.


\subsection{Základní přístupy k~doporučování}
\label{subsec:zakladni_pristupy}

Podle standardní oborové taxonomie, kterou definoval Burke \cite{Burke2007} a~ve svém přehledu shrnují Ricci a~kol. \cite{RicciRokachShapira2022}, lze doporučovací systémy rozdělit do několika kategorií podle toho, s~jakými daty pracují. Z~hlediska historického vývoje i~současné praxe se nejčastěji setkáváme se třemi hlavními přístupy:

\begin{itemize}
    \item \textbf{Kolaborativní filtrování (CF):} Tento přístup staví výhradně na historii interakcí mezi uživateli a~položkami. Předpokládá, že uživatelé se stejným chováním v~minulosti (např. navštěvující stejné typy restaurací) budou mít podobný vkus i~v~budoucnu \cite{RicciRokachShapira2022}. Ačkoliv jde o~vysoce účinnou metodu, naráží na kritický problém tzv. \textit{studeného startu} (cold start). Ten nastává v~okamžiku, kdy do systému přibude zcela nový uživatel či nová položka bez jakékoliv historie interakcí. Bez historických dat nedokáže CF najít relevantní podobnosti a~selhává.
    \item \textbf{Obsahové filtrování (CBF):} Tento model naopak porovnává vnitřní vlastnosti (metadata) samotných položek s~uživatelským profilem \cite{Burke2007}. V~kontextu doporučování bodů zájmu (POI) může jít o~textové popisy lokací, přiřazené tagy či kategorie (např. italská kuchyně, bezbariérový přístup). Protože k~hodnocení nevyžaduje data od ostatních uživatelů, je CBF přirozeně odolné vůči problému studeného startu u~nových položek. Hrozí u~něj ovšem tzv. \textit{přespecifikování} (overspecialization), kdy systém uživateli nabízí pouze položky velmi podobné těm, které již zná, a~ztrácí schopnost objevovat pro něj nové koncepty.
    \item \textbf{Hybridní přístupy:} Tyto metody cíleně kombinují výhody obou předchozích světů \cite{Burke2007, RicciRokachShapira2022}. Využitím obsahových rysů k~obohacení řídkých kolaborativních dat dokáží elegantně překonat úskalí studeného startu, a~zároveň si díky kolaborativní složce zachovávají schopnost objevovat pro uživatele nové a~často nečekané okruhy zájmů.
\end{itemize}

Abychom v~následujících kapitolách mohli přesně popsat, jakým způsobem náš navrhovaný systém tyto přístupy kombinuje, je nezbytné si nejprve ukázat, jak jsou v~doporučovacích systémech data a~cíle reprezentovány matematicky. Ačkoliv se totiž výše uvedené metody liší v~tom, jaké informace zpracovávají, všechny sdílejí stejný fundamentální úkol: odhadnout neznámou míru preference konkrétního uživatele k~dané položce. 

Nechť $U = \{u_1, u_2, \dots, u_{|U|}\}$ je množina všech uživatelů a~$I = \{i_1, i_2, \dots, i_{|I|}\}$ je množina všech dostupných položek. Historie pozorovaných interakcí je v~základním pojetí zachycena pomocí matice $\mathbf{X} \in \mathbb{R}^{|U| \times |I|}$, kde prvek $x_{u,i}$ nabývá konkrétní hodnoty v~případě, že mezi uživatelem a~položkou proběhla interakce (např. hodnota 1 pro navštívené místo, 0 jinak). 

Cílem doporučovacího systému je naučit se predikční funkci $f$, která pro libovolný pár uživatel-položka odhadne chybějící relevanci (tzv. skóre) jako $\hat{x}_{u,i} = f(u, i \,|\, \Theta)$, kde $\Theta$ představuje natrénované parametry modelu. Uvedené tři přístupy se pak formálně liší právě v~tom, z~čeho se tato funkce učí. Zatímco u~čistého kolaborativního filtrování se parametry $\Theta$ optimalizují výhradně ze vztahů uvnitř matice interakcí $\mathbf{X}$, v~případě obsahových a~hybridních metod zohledňuje funkce $f$ i~dodatečné vektory vlastností (metadat) samotných položek. Právě to umožňuje modelu matematicky vyčíslit skóre $\hat{x}_{u,i}$ i~pro položku s~prázdnou historií interakcí, čímž je formálně překonán onen zmíněný problém studeného startu.

\subsection{Specifika doporučování bodů zájmu (POI)}
\label{subsec:specifika_poi}

Doporučování bodů zájmu (\textit{Points of Interest}, POI) představuje specifickou podoblast doporučovacích úloh. Na rozdíl od tradičních platforem pro elektronický obchod (kde si uživatel produkt pouze koupí a~nechá doručit) je výběr volnočasových aktivit silně ovlivněn fyzickým a~situačním kontextem. Mezi nejvýznamnější omezující faktory patří:

\begin{itemize}
    \item \textbf{Geografická poloha:} Doporučení je pro uživatele relevantní pouze tehdy, nachází-li se místo v~jeho fyzické dostupnosti a~vzdálenosti, kterou je ochoten překonat.
    \item \textbf{Časové a~situační podmínky:} Dostupnost lokace je vázána na otevírací dobu, přičemž vhodnost doporučení může dramaticky měnit i~počasí nebo účel cesty (např. oběd s~rodinou vs. večerní posezení s~přáteli).
\end{itemize}

V~současném výzkumu se k~těmto lokačním datům přistupuje dvěma hlavními způsoby. Prvním z~nich je modelování problému jako uspořádané sekvence návštěv (tzv. \textit{Next-POI recommendation}). Zde se předpokládá, že interakce uživatele $u$ tvoří časovou řadu $S_u = \{(i_1, t_1), (i_2, t_2), \dots, (i_n, t_n)\}$, a~systém se snaží odhadnout bezprostředně následující krok s~ohledem na časové rozestupy a~překonanou vzdálenost \cite{Zeng2025}. 

Ačkoliv jsou sekvenční modely vysoce efektivní pro predikci tras v~reálném čase, pro hlubší porozumění zájmům uživatele (např. při dlouhodobém plánování dovolené) bývá vhodnější \textbf{druhý přístup, založený na ne-sekvenčním modelování preferencí} \cite{Ye2010, Lian2014_GeoMF}. Ten abstrahuje od striktní časové návaznosti a~nahlíží na problém skrze celkovou matici historických interakcí (typicky s~využitím maticové faktorizace či kolaborativního filtrování). Systém tak dokáže vytvořit komplexní profil uživatele (např. "má rád italskou kuchyni a~historické památky"), který je stabilnější v~čase. Tento přístup je ústředním prvkem i~pro námi navrhovanou architekturu, u~níž jsou pro zachycení zájmů primárně využívány metody hlubokého kolaborativního filtrování nad statickou maticí interakcí.

Bez ohledu na zvolenou metodu modelování je však v~doméně POI důležitá \textbf{vysvětlitelnost a~transparentnost}. Zatímco u~digitálního obsahu (např. u~filmů či hudby) lze špatné doporučení snadno přeskočit, návštěva fyzické lokace vyžaduje od uživatele investici času a~často i~peněz. Proto je pro budování důvěry nezbytné, aby systém dokázal svá rozhodnutí přesvědčivě a~srozumitelně zdůvodnit. Lokační data navíc představují specifický problém v~tom, že v~sobě přirozeně mísí (tzv. \textit{entanglují}) zcela odlišné motivace. Jak upozorňují Zeng a~kol. \cite{Zeng2025}, check-in chování je silně ovlivněno prostorovými a~časovými faktory, které se v~datech složitě prolínají se skutečnými sémantickými zájmy uživatele. Běžné modely tyto vlivy spojují do jedné černé skříňky, čímž ztrácejí expresivitu. 

Uživatel přitom potřebuje přesně chápat, zda mu systém doporučuje konkrétní restauraci kvůli jejím vnitřním vlastnostem (výborné hodnocení,~sémantická shoda s~jeho chutěmi), nebo čistě proto, že se zrovna nachází v~bezprostřední blízkosti. Právě nutnost tyto složité motivace rozplést (tzv. \textit{disentanglement}) do srozumitelných pojmů \cite{Zeng2025} nás přivádí k~využití pokročilých metod skrytých reprezentací, které jsou popsány v~následující sekci.

\section{Latentní reprezentace a~architektura autoenkodérů}
\label{sec:uceni_reprezentaci}

Klasické metody kolaborativního (CF) a~obsahového filtrování (CBF) sice dosahují dobrých výsledků, avšak jejich vnitřní (latentní) reprezentace jsou často složité a~pro uživatele i~vývojáře obtížně interpretovatelné. Moderní přístupy proto stále častěji využívají modely hlubokého učení, zejména rodinu autoenkodérů, jejichž hlavní výhodou je schopnost zachytit i~velmi komplexní a~nelineární vztahy mezi uživateli a~položkami. 

\subsection{Autoenkodéry pro kolaborativní filtrování (CFAE)}
\label{subsec:cfae}

Princip těchto modelů spočívá v~tom, že se autoenkodér učí komprimovat řídká interakční data do zmenšeného latentního prostoru a~následně je z~něj s~co nejmenší chybou zrekonstruovat. V~architektuře CFAE (\textit{Collaborative-Filtering Autoencoders}, zastoupené například modely MultVAE nebo lineární sítí ELSA) je každý uživatel reprezentován vektorem $x_u \in \mathbb{R}^{|I|}$. Enkodér $E_c$ komprimuje tento dlouhý a~řídký vektor do hustého (\textit{dense}) latentního prostoru o~nízké dimenzionalitě $d$:

\begin{equation}
y_u = E_c(x_u) = \sigma(W_{enc} x_u + b_{enc})
\end{equation}

Dekodér $D_c$ se následně snaží tento zkomprimovaný vektor $y_u$ zrekonstruovat a~predikovat skóre pro chybějící položky jako $\hat{x}_u = D_c(y_u)$. 

Přestože tyto modely nabízejí vynikající prediktivní přesnost, jejich interpretace je obtížná. Vytvářejí totiž tzv. propletené (\textit{entangled}) reprezentace, kde jedna dimenze ovlivňuje mnoho zcela nesouvisejících vlastností položky (např. zároveň žánr restaurace i~její lokaci). To z~těchto modelů dělá neprůhledné černé skříňky (\textit{black-boxes}) \cite{Spisak2024}. Vzhledem k~potřebě vysvětlitelnosti v~doméně lokačních dat je proto nutné tyto reprezentace dodatečně rozplést.

\subsection{Řídké autoenkodéry (SAE)}
\label{subsec:sae}

Problém černých skříněk a~propletených vektorů lze řešit architekturou řídkých autoenkodérů (\textit{Sparse Autoencoders}, SAE). Ty představují funkční rozšíření klasických autoenkodérů, které do trénovacího procesu uměle zavádí přísný požadavek na řídkost (\textit{sparsity}) v~latentní vrstvě. Na rozdíl od běžného autoenkodéru se SAE snaží omezit počet současně aktivních neuronů. Toho lze dosáhnout například aktivační funkcí ReLU v~kombinaci s~$L1$ regularizací, nebo použitím mechanismu \textit{Top-K} výběru, který propustí pouze $K$ nejsilnějších signálů. Tímto omezením se model učí kompaktnější a~sémanticky smysluplnější reprezentace, a~vyhýbá se tak pouhému zapamatování si dat.

Z~matematického hlediska přebírá modul SAE onen hustý a~neprůhledný vektor $y_u \in \mathbb{R}^d$ ze základního modelu CFAE a~mapuje ho do podstatně širšího prostoru $\mathbb{R}^k$ (kde $k \gg d$) \cite{Spisak2024}:

\begin{equation}
z_u = E_s(y_u) = \phi(W_s y_u + b_s)
\end{equation}

V~této rovnici tvoří $z_u$ hledanou rozpletenou reprezentaci a~funkce $\phi$ vynucuje požadovanou řídkost. Následná celková ztrátová funkce pro trénování tohoto modulu kombinuje chybu rekonstrukce původního hustého vektoru a~penalizaci vynucující řídkost (čímž brání tvorbě příliš hustých reprezentací):

\begin{equation}
\mathcal{L}_{SAE} = \frac{1}{|U|} \sum_{u \in U} \left( ||y_u - \hat{y}_u||_2^2 \right) + \lambda ||z_u||_1
\end{equation}

\subsection{Interpretace a~ovládací prvky (Knobs)}
\label{subsec:interpretace_knobs}

Omezením aktivační kapacity prostřednictvím SAE získává výsledná reprezentace $z_u$ unikátní schopnost -- její jednotlivé dimenze začínají odpovídat samostatným a~izolovaným významovým faktorům v~datech. Jak upozorňuje Wang a~kol. \cite{Wang2021}, klasické modely nedokážou tyto faktory přirozeně oddělit, což silně limituje jejich praktickou využitelnost pro pokročilé uživatelské interakce.

Zavedením řídkosti však autoenkodér začne fungovat jako sémantický filtr, který složitě propletená (entanglovaná) data pomyslně rozčesává do struktury nezávislých vláken významu \cite{Spisak2024}. Pokud například konkrétní dimenze kóduje „zájem o~italskou kuchyni“, je možné ji uživateli takto přímo prezentovat. V~odborné literatuře se pro tyto rozpletené rysy vžilo označení \textit{ovládací prvky} či \textit{knobs} \cite{Spisak2024, ONeill2024}. Jejich zásadní výhodou je, že je lze přímo napojit na uživatelské rozhraní. Tím se otevírá cesta k~tomu, aby systém nejen dokázal svá doporučení zpětně vysvětlit, ale také aby umožnil samotnému uživateli do procesu aktivně zasahovat.

\section{Vysvětlitelnost a~interaktivní řízení (Steering)}
\label{sec:vysvetlitelnost_steering}

Jakmile je složitý latentní prostor modelu rozpleten do srozumitelných sémantických prvků, lze tyto reprezentace využít k~obousměrné komunikaci mezi uživatelem a~systémem. Směrem od systému k~uživateli slouží k~poskytování transparentních vysvětlení, zatímco směrem od uživatele k~systému otevírají prostor pro aktivní řízení (steering).

\subsection{Vysvětlitelnost doporučení a~její hodnocení}
\label{subsec:vysvetlitelnost}

Vysvětlitelnost (\textit{Explainability}) znamená schopnost modelu srozumitelně sdělit uživateli důvody, proč mu byla konkrétní položka navržena. Využití architektury SAE umožňuje dosáhnout tzv. \textit{model-based} (neboli \textit{white-box}) vysvětlitelnosti, kdy je vysvětlení integrální součástí algoritmu a~vychází přímo z~jeho vnitřního matematického stavu. 

Při využití řídkého autoenkodéru lze totiž celkové skóre doporučení formalizovat zjištěním příspěvků jednotlivých složek řídkého vektoru $z_u$. Pokud budeme předpokládat, že je dekodér lineární vrstvou $V$, můžeme predikované skóre $\hat{x}_{u,i}$ pro položku $i$ rozepsat jako sumu příspěvků jednotlivých aktivních neuronů:

\begin{equation}
\hat{x}_{u,i} \approx \sum_{j=1}^k (V_i \cdot W_{s,j}^T) z_{u,j}
\end{equation}

Díky této dekompozici dokáže systém uživateli přesně odůvodnit, jaký specifický sémantický rys k~finálnímu doporučení přispěl nejvíce \cite{Spisak2024}. 

Hodnocení kvality takto vygenerovaných vysvětlení představuje v~odborné literatuře komplexní výzvu a~rozlišují se při něm dva hlavní směry. Z~objektivního hlediska se měří tzv. \textbf{fidelity} (věrnost), která matematicky vyjadřuje, nakolik prezentované vysvětlení skutečně odpovídá skutečnému vnitřnímu rozhodovacímu procesu modelu. Druhým, avšak pro praktické nasazení často důležitějším přístupem, jsou \textbf{uživatelské studie}. Ty hodnotí vysvětlení z~pohledu koncového uživatele a~zaměřují se na metriky, jakými jsou vnímaná srozumitelnost, přesvědčivost (persuasiveness) či schopnost budovat důvěru (trust) \cite{Tintarev2012}. Vysvětlení totiž nemusí být jen technicky přesné, ale musí především rezonovat s~mentálním modelem uživatele.

\subsection{Steering: Definice a~pohledy na ovladatelnost}
\label{subsec:steering}

Pokud uživatel díky vysvětlení pochopí, jak model uvažuje, přirozeně vyžaduje možnost toto uvažování korigovat. Pojem \textit{steering} (řízení či kormidlování) se v~literatuře o~doporučovacích systémech objevuje stále častěji, avšak nenabývá jediné ustálené definice. V~současném výzkumu lze identifikovat tři odlišné pohledy na to, kdo a~jakým způsobem doporučování řídí:

\begin{enumerate}
    \item \textbf{Algoritmický steering platforem:} Z~pohledu ekonomiky a~návrhu platforem (často označováno jako \textit{self-preferencing}) představuje steering situaci, kdy samotný algoritmus vychyluje doporučení ve prospěch určitých položek na úkor čisté uživatelské preference \cite{Chen2024}. Typickým příkladem je, když platforma (např. Amazon) algoritmicky "směruje" uživatele ke svým vlastním produktům místo k~relevantnějším produktům třetích stran \cite{Chen2024}. Zde je tedy uživatel objektem řízení.
    \item \textbf{Strategický steering ze strany uživatele:} Zcela opačný extrém definují Colella a~kol. \cite{Colella2020}. Upozorňují, že uživatelé nejsou pouze pasivními poskytovateli dat, ale inteligentními strategickými aktéry. Pokud například uživatel chce dostávat více doporučení na fantasy filmy, může začít záměrně (a~nepravdivě) udělovat vysoká hodnocení položkám, o~kterých si myslí, že tento žánr v~očích algoritmu posílí. Uživatel tak pomocí falešné zpětné vazby strategicky manipuluje (steeruje) systém, aby dosáhl svého osobního cíle \cite{Colella2020}.
    \item \textbf{Explicitní interaktivní řízení (User-controllability):} Prostřední a~z~hlediska této práce nejrelevantnější přístup chápe steering jako přímou, transparentní a~podporovanou funkci uživatelského rozhraní (HCI). Výzkumy jako systém TasteWeights \cite{Tasteweights2012} či studie v~oblasti e-learningu \cite{Ooge2023} prokazují, že pokud dáme uživatelům explicitní vizuální nástroje (např. posuvníky) k~úpravě vah jejich preferencí, dramaticky to zvyšuje jejich důvěru a~pocit kontroly (agency) nad systémem.
\end{enumerate}

V~námi uvažované architektuře se steering realizuje právě třetím jmenovaným způsobem, a~to přímou modifikací hodnot uživatelského vektoru v~latentním prostoru SAE. Nechť uživatel prostřednictvím posuvníku v~rozhraní vyjádří záměr posílit svůj zájem o~konkrétní sémantický rys (např. „hledám více muzeí“), jemuž odpovídá neuron $n$. Systém zkonstruuje nový, modifikovaný preferenční vektor $z'_u$ pomocí afinní interpolace:

\begin{equation}
z'_u = \alpha z_u + (1 - \alpha) e_n
\end{equation}

kde $e_n \in \mathbb{R}^k$ je tzv. one-hot vektor s~hodnotou 1 na pozici žádaného neuronu $n$ a~0 jinde. Parametr $\alpha \in [1]$ určuje sílu zachování původního profilu. Čím více se $\alpha$ blíží k~nule, tím agresivněji systém přepíše uživatelovu historickou stopu ve prospěch nově vyžádané preference. Takto upravený vektor je následně v~reálném čase propuštěn skrze dekodér, čímž dojde k~okamžité aktualizaci doporučených položek.

\subsection{Etické, právní a~kognitivní aspekty interakce}
\label{subsec:etika_pravo_kognice}

Poskytnutí takto silné ovladatelnosti otevírá důležité otázky z~oblasti psychologie rozhodování. Zviditelnění dopadů řízení (steeringu) pomáhá předcházet tzv. nadměrnému či nedostatečnému spoléhání se na systém (\textit{over- and under-reliance}) \cite{Tintarev2012}. Uživatel, který vidí, z~čeho doporučení vychází, a~může ho upravit, dokáže snáze ignorovat irelevantní návrhy. Na straně druhé může přílišná kontrola vést k~nežádoucí akceleraci tvorby informačních bublin (\textit{filter bubbles}). Jakmile si uživatel systém extrémně přizpůsobí pouze svým aktuálním a~ověřeným zájmům, ztrácí se prvek objevování nového a~nečekaného obsahu (tzv. \textit{serendipity}). 

Kromě kognitivních aspektů je však nutné zohlednit i~širší \textbf{etické implikace}. Jak upozorňují studie zaměřené na férovost (\textit{fairness}) v~doporučovacích systémech \cite{Ekstrand2022_Fairness}, algoritmy mohou snadno zrcadlit či dokonce zesilovat strukturální předsudky a~tzv. \textit{popularity bias} (neúměrné upřednostňování již populárních položek na úkor těch méně známých). Další etický problém nastává, pokud má poskytovatel platformy duální roli -- tedy funguje jako informační zprostředkovatel i~jako přímý prodejce. V~takových případech může docházet k~netransparentnímu algoritmickému steeringu ze strany platformy (tzv. \textit{self-preferencing}), kdy systém skrytě směruje uživatele k~vlastním či ziskovějším položkám \cite{Chen2024}. Transparentní a~uživatelem explicitně ovladatelný profil tak představuje silný nástroj, jak tuto informační asymetrii a~riziko algoritmické manipulace omezit.

Tyto etické a~psychologické výzvy se navíc v~posledních letech promítají přímo do \textbf{legislativní roviny}. V~kontextu Evropské unie zavedlo Obecné nařízení o~ochraně osobních údajů (GDPR) přísnější požadavky na transparentnost automatizovaného rozhodování, které se v~odborné literatuře o~umělé inteligenci často interpretují jako tzv. \textit{právo na vysvětlení} (right to explanation) \cite{Goodman2017_GDPR, Musto2022_Semantics}. Od systémů, které významně ovlivňují uživatelské chování a~pracují s~osobními preferencemi, se tak dnes již nevyžaduje pouze vysoká prediktivní přesnost, ale i~právně podložená schopnost svá rozhodnutí srozumitelně zdůvodnit. Návrh moderních interaktivních doporučovacích systémů by proto měl vždy balancovat mezi uspokojením přesných uživatelských požadavků, zachováním dostatečné rozmanitosti a~dodržováním oborových i~právních standardů transparentnosti.

\section{Shrnutí}
\label{sec:shrnuti_teorie}

Tato kapitola položila teoretické a~formální základy pro návrh vysvětlitelného a~ovladatelného doporučovacího systému. V~první části jsme představili klasické přístupy k~doporučování a~zasadili je do specifického kontextu bodů zájmu (POI). Ukázalo se, že v~této doméně, kde každé doporučení vyžaduje fyzickou a~časovou investici uživatele, je transparentnost systému a~schopnost rozplést složité lokační motivace naprosto klíčová.

Následně jsme detailně popsali přechod od řídkých matic interakcí k~latentním reprezentacím pomocí hlubokého učení. Zatímco moderní autoenkodéry pro kolaborativní filtrování (CFAE) dosahují špičkové prediktivní přesnosti, vytvářejí neprůhledné, tzv. propletené černé skříňky. Bylo formálně ukázáno, jak architektura řídkých autoenkodérů (SAE) tvoří funkční most mezi těmito přesnými modely a~požadavkem na srozumitelnost. Vynucením řídkosti dokáže SAE tyto nesrozumitelné vektory převést na izolované, monosémantické rysy -- tzv. ovládací prvky (knobs).

V~poslední části jsme vysvětlili, že jakmile je latentní prostor takto rozpleten, otevírá se cesta k~pokročilé obousměrné interakci. Směrem k~uživateli umožňuje systém generovat exaktní a~věrná vysvětlení přímo na základě vnitřního stavu sítě (\textit{model-based explainability}), zatímco směrem k~systému poskytuje uživateli možnost aktivně a~v~reálném čase řídit svůj profil (steering). Byla přitom zdůrazněna i~nutnost brát v~potaz psychologické a~kognitivní aspekty, jako je budování důvěry na jedné straně a~riziko vzniku informačních bublin na straně druhé.

Z~tohoto teoretického ukotvení budeme pevně vycházet v~následujících částech práce. Zatímco následující kapitola kriticky zmapuje, jakým způsobem k~těmto jednotlivým výzvám dosud přistupovala a~přistupuje světová odborná literatura, třetí kapitola tyto poznatky následně přetaví do samotného architektonického návrhu našeho řešení.